AlphaFold DB covers more than 200 million proteins, and the rule ``\plddt $<$ 50 means
intrinsically disordered'' has become an everyday working heuristic for trimming models and
annotating the dark proteome. The rule has never been validated at the operating point at
which it is used: the literature reports ROC-AUC, a threshold-free summary that cannot
adjudicate a claim about a specific cutoff. We audit the rule directly on the CAID round-2 and
round-3 blind assessments, across both official reference conventions (\dpdb and \dnox), with
bootstrap confidence intervals over targets and coverage-matched target sets, against a panel
of 37--56 dedicated disorder predictors. We find that \plddtrule has sensitivity 0.60--0.63
(0.63--0.65 restricted to regions longer than 30 residues), not above 0.8: not one of 1000
bootstrap resamples reaches 0.80 in any of four settings. Its specificity is 0.98 under \dpdb
and 0.67--0.76 under \dnox --- a 30-point swing produced entirely by how the benchmark defines
a negative, not by AlphaFold. Once every method is threshold-calibrated, none of 174
method\,$\times$\,dataset comparisons beats \plddt on short ($<$20 aa) segments, and the median
dedicated predictor is 1.5--8.7 points \emph{worse}; the widely assumed 10--15\% advantage is
exactly reproducible at a shared naive 0.5 cut and vanishes after calibration. What survives is
the mechanism: reading solvent accessibility out of the \emph{same} AlphaFold2 coordinates is
significantly better on every reference set ($\Delta$AUC $+0.015$ to $+0.056$, all $p \le 0.02$),
and the disorder \plddt misses is confidently modelled and solvent-exposed.
\plddtrule is a high-precision, low-recall detector; users who need recall should raise the
cutoff to $\approx$70 and prefer the geometric readout of the same model.
